% GENERATED FILE. Edit canonical lessons, then run npm run generate:books.
% canonical-source-hash: fnv1a64:4582787eafb88ec1
% canonical-lessons: RU-C161-anketa, RU-C161-konvert, RU-C161-marka, RU-C161-otkrytka, RU-C161-posylka

\chapter{Form, Envelope, Postage stamp, Postcard, Parcel}
\label{ch:ru-form-envelope-postage-stamp-postcard-parcel}
\hlchaptermodality{161}

\begin{chapteropening}
\textbf{By the end of this chapter:} \emph{I can say a form, an envelope, a postage stamp, a postcard and a parcel.}

Complete the last lesson of chapter 161: I can say a form, an envelope, a postage stamp, a postcard and a parcel.
\end{chapteropening}

\section[\texorpdfstring{\ru{анкета} (ankéta)}{анкета (ankéta)}]{\ru{анкета} (ankéta) — a form, a questionnaire}

\label{lesson:RU-C161-anketa}



\begin{quote}
Before the new one: say the Russian for a drawer; a crate, then the Russian for a light bulb.
\end{quote}

\subsection*{You'll want to know: \ru{анкета}}
\textbf{\ru{анкета}} (\emph{ankéta}) — \textquotedblleft{}a form, a questionnaire\textquotedblright{}.

\begin{grammarlens}[title={What you've built}]
One more word for this chapter.
\end{grammarlens}

\subsection*{Your turn}
\begin{itemize}
\raggedright
  \item \emph{Say it:} \emph{ankéta}
  \item \emph{Say it:} \emph{ankéta}, once more
  \item \emph{Say it:} say \emph{ankéta}
  \item \emph{Recall:} say the Russian for a drawer; a crate, then the Russian for a light bulb, then say \emph{ankéta} again
\end{itemize}

\begin{tcolorbox}[breakable,colback=teal!4,colframe=teal!35!black,title={Before you move on}]
What does \ru{анкета} mean? (\textquotedblleft{}A form, a questionnaire\textquotedblright{}.) Say it once more.
\end{tcolorbox}
\section[\texorpdfstring{\ru{конверт} (konvért)}{конверт (konvért)}]{\ru{конверт} (konvért) — an envelope}

\label{lesson:RU-C161-konvert}



\begin{quote}
Before the new one: say the Russian for a light bulb, then the Russian for a form.
\end{quote}

\subsection*{You'll want to know: \ru{конверт}}
\textbf{\ru{конверт}} (\emph{konvért}) — \textquotedblleft{}an envelope\textquotedblright{}.

\begin{grammarlens}[title={What you've built}]
One more word for this chapter.
\end{grammarlens}

\subsection*{Your turn}
\begin{itemize}
\raggedright
  \item \emph{Say it:} \emph{konvért}
  \item \emph{Say it:} \emph{konvért}, once more
  \item \emph{Say it:} say \emph{konvért}
  \item \emph{Recall:} say the Russian for a light bulb, then the Russian for a form, then say \emph{konvért} again
\end{itemize}

\begin{tcolorbox}[breakable,colback=teal!4,colframe=teal!35!black,title={Before you move on}]
What does \ru{конверт} mean? (\textquotedblleft{}An envelope\textquotedblright{}.) Say it once more.
\end{tcolorbox}
\section[\texorpdfstring{\ru{марка} (márka)}{марка (márka)}]{\ru{марка} (márka) — a postage stamp; a brand}

\label{lesson:RU-C161-marka}



\begin{quote}
Before the new one: say the Russian for a form, then the Russian for an envelope.
\end{quote}

\subsection*{You'll want to know: \ru{марка}}
\textbf{\ru{марка}} (\emph{márka}) — \textquotedblleft{}a postage stamp; a brand\textquotedblright{}.

\begin{grammarlens}[title={What you've built}]
One more word for this chapter.
\end{grammarlens}

\subsection*{Your turn}
\begin{itemize}
\raggedright
  \item \emph{Say it:} \emph{márka}
  \item \emph{Say it:} \emph{márka}, once more
  \item \emph{Say it:} say \emph{márka}
  \item \emph{Recall:} say the Russian for a form, then the Russian for an envelope, then say \emph{márka} again
\end{itemize}

\begin{tcolorbox}[breakable,colback=teal!4,colframe=teal!35!black,title={Before you move on}]
What does \ru{марка} mean? (\textquotedblleft{}A postage stamp; a brand\textquotedblright{}.) Say it once more.
\end{tcolorbox}
\section[\texorpdfstring{\ru{открытка} (otkrýtka)}{открытка (otkrýtka)}]{\ru{открытка} (otkrýtka) — a postcard}

\label{lesson:RU-C161-otkrytka}



\begin{quote}
Before the new one: say the Russian for an envelope, then the Russian for a postage stamp; a brand.
\end{quote}

\subsection*{You'll want to know: \ru{открытка}}
\textbf{\ru{открытка}} (\emph{otkrýtka}) — \textquotedblleft{}a postcard\textquotedblright{}.

\begin{grammarlens}[title={What you've built}]
One more word for this chapter.
\end{grammarlens}

\subsection*{Your turn}
\begin{itemize}
\raggedright
  \item \emph{Say it:} \emph{otkrýtka}
  \item \emph{Say it:} \emph{otkrýtka}, once more
  \item \emph{Say it:} say \emph{otkrýtka}
  \item \emph{Recall:} say the Russian for an envelope, then the Russian for a postage stamp; a brand, then say \emph{otkrýtka} again
\end{itemize}

\begin{tcolorbox}[breakable,colback=teal!4,colframe=teal!35!black,title={Before you move on}]
What does \ru{открытка} mean? (\textquotedblleft{}A postcard\textquotedblright{}.) Say it once more.
\end{tcolorbox}
\section[\texorpdfstring{\ru{посылка} (posýlka)}{посылка (posýlka)}]{\ru{посылка} (posýlka) — a parcel}

\label{lesson:RU-C161-posylka}



\begin{quote}
Before the new one: say the Russian for a postage stamp; a brand, then the Russian for a postcard.
\end{quote}

\subsection*{You'll want to know: \ru{посылка}}
\textbf{\ru{посылка}} (\emph{posýlka}) — \textquotedblleft{}a parcel\textquotedblright{}.

\begin{grammarlens}[title={What you've built}]
One more word for this chapter.
\end{grammarlens}

\subsection*{Your turn}
\begin{itemize}
\raggedright
  \item \emph{Say it:} \emph{posýlka}
  \item \emph{Say it:} \emph{posýlka}, once more
  \item \emph{Say it:} say \emph{posýlka}
  \item \emph{Recall:} say the Russian for a postage stamp; a brand, then the Russian for a postcard, then say \emph{posýlka} again
\end{itemize}

\begin{tcolorbox}[breakable,colback=teal!4,colframe=teal!35!black,title={Before you move on}]
What does \ru{посылка} mean? (\textquotedblleft{}A parcel\textquotedblright{}.) Say it once more.
\end{tcolorbox}
